% Default preamble for typeset notes: classic LaTeX article style, in the style of a standard homework write-up.
% (Computer Modern, bold title, italic introduction, unindented justified paragraphs, numbered equations,
%  nested lists 1. / (a) / i. / A., boxed tables with italic "Table n:" captions.)
\usepackage[letterpaper,margin=0.5in]{geometry}
\usepackage{amsmath,amssymb,amsthm,mathtools}
\usepackage{lmodern}
\usepackage[T1]{fontenc}
\usepackage{enumitem}
\usepackage{array}
\usepackage[labelsep=colon,textfont=it]{caption}
\usepackage{braket}
\usepackage[hidelinks]{hyperref}
\usepackage{microtype}

\setlength{\parindent}{0pt}
\setlength{\parskip}{0.6em}
\setlist[enumerate,1]{label=\arabic*.,leftmargin=2.2em}
\setlist[enumerate,2]{label=(\alph*),leftmargin=2em}
\setlist[enumerate,3]{label=\roman*.,leftmargin=2em}
\setlist[enumerate,4]{label=\Alph*.,leftmargin=1.6em,nosep}

\newtheorem{theorem}{Theorem}[section]
\newtheorem{lemma}[theorem]{Lemma}
\newtheorem{proposition}[theorem]{Proposition}
\newtheorem{corollary}[theorem]{Corollary}
\theoremstyle{definition}
\newtheorem{definition}[theorem]{Definition}
\newtheorem{example}[theorem]{Example}
\theoremstyle{remark}
\newtheorem*{remark}{Remark}

\newcommand{\R}{\mathbb{R}}
\newcommand{\C}{\mathbb{C}}
\newcommand{\Z}{\mathbb{Z}}
\newcommand{\N}{\mathbb{N}}
\newcommand{\Q}{\mathbb{Q}}
\renewcommand{\Re}{\mathcal{R}}
\renewcommand{\Im}{\mathcal{I}}

% Title block: bold title, then an italic introductory paragraph.
\newcommand{\notetitle}{Notes}
\newcommand{\maketitleblock}[1]{{\noindent\Large\bfseries #1\par}\vspace{0.6em}}
\newenvironment{intro}{\itshape}{\par}
